\section{Spectral input and tree identities}\label{sec:preliminaries}

Throughout the proof, \(d\ge3\) and \(|E|<2\sqrt{d-1}\) are fixed.
All constants may depend on these two parameters.
This section records the spectral information used below and the tree
identities that determine the covariance of the eigenvector entries.
The spectral estimates are established results; the subsequent entropy
and switching arguments are proved in this paper.

\subsection{Conventions and microscopic coordinates}

Let
\[
 \rho_d(t)=
 \frac{d\sqrt{4(d-1)-t^2}}{2\pi(d^2-t^2)}
 \mathbf 1_{\{|t|<2\sqrt{d-1}\}},
 \qquad
 m_d(\zeta)=\int_{\R}\frac{\rho_d(t)}{t-\zeta}\dd t
 \quad(\Im\zeta>0).
\]
Set
\[
 \rho=\rho_d(E),\qquad p=\frac Ed,\qquad
 a=\frac{d}{d-1}(1-p^2),\qquad
 m_*=m_d(E+\mathrm i0)=m_{\rm r}+\mathrm i\pi\rho.
\]
In particular, \(\rho>0\), \(a>0\), and \(|p|<1\).
For the adjacency eigenvalues and a graph-determined real orthogonal
eigenbasis, use the normalization
\[
 Au_i=\lambda_i u_i,\qquad
 \langle u_i,u_j\rangle=n\delta_{ij}.
\]
An arbitrary measurable choice inside each eigenspace suffices.
Define
\begin{equation}\label{eq:micro-coordinates}
 x_i=n(\lambda_i-E),\qquad
 \nu_n=\sum_{i=1}^n\delta_{x_i},\qquad
 m_n(z)=\sum_{i=1}^n\frac1{x_i-z}
       =\frac1n\Tr R_n(z),
 \quad
 R_n(z)=(A-E-z/n)^{-1}.
\end{equation}
The remaining density factor for unfolding is \(\rho\); the final
point coordinates are obtained by multiplying the \(x_i\) by \(\rho\).
We write \(\C_+=\{z\in\C:\Im z>0\}\) for the upper half-plane.

For \(k\ge1\), let
\begin{equation}\label{eq:channel-matrices}
 J_k=\begin{pmatrix}0&I_k\\ I_k&0\end{pmatrix},
 \qquad C_k=I_{2k}+pJ_k,\qquad
 K_k(m)=mC_k+d^{-1}J_k.
\end{equation}
The order of the \(2k\) coordinates will always be the starting endpoints
of \(k\) oriented edges, followed by their ending endpoints.
The matrix \(C_k\) is positive definite.

Unsubscripted finite-dimensional matrix norms are operator norms;
\(\|\cdot\|_{\HS}\) denotes the Hilbert--Schmidt norm.
The imaginary part of a matrix is its Hermitian imaginary part.
For the complex symmetric matrices occurring here it also agrees with
entrywise imaginary part.
All tensor norms are induced by the stated Euclidean inner products.
The notation \(n^{o(1)}\), when used on high-probability events, denotes a
deterministic subpolynomial upper bound.
We write \(O_{\PP}(b_r)\) for a random quantity whose ratio to \(b_r\) is
bounded in probability; \(o_{\PP}(b_r)\) means that the ratio tends
to zero in probability.

\subsection{The simple regular-graph model}

We use the classical pairing model for regular graphs
\cite{BollobasPairing}. The special counting facts needed here have
the following direct proof.

\begin{lemma}\label{lem:model}
In the configuration model with \(d\) stubs at each of \(n\) labelled
vertices, conditioning the matching to give a simple graph produces the
uniform simple labelled \(d\)-regular graph. Moreover,
\[
 \PP(\text{the matching gives a simple graph})
       \longrightarrow \exp\!\left(-\frac{d^2-1}{4}\right)>0
\]
along all admissible sizes.
In the simple model, a specified vertex belongs to a triangle with
probability \(O(n^{-1})\).
For any fixed \(k\), independently sampled uniform oriented edges have
distinct endpoints and induce exactly a matching with probability
\(1-O_k(n^{-1})\).
\end{lemma}

\begin{proof}
Every simple graph has the same number \((d!)^n\) of stub realizations,
which proves the conditional uniformity.
A prescribed loop is one pair of stubs at a vertex. A prescribed
double edge is an unordered pair of pairs between two distinct vertices,
using distinct stubs. Let \(Z\) count all occurring prescribed defects
of these two kinds.
The expected counts tend to
\[
 \frac{d-1}{2}\quad\hbox{and}\quad\frac{(d-1)^2}{4},
\]
respectively.
Indeed, a prescribed family of \(e\) disjoint pairs occurs with
probability
\[
 \frac{1}{(nd-1)(nd-3)\cdots(nd-2e+1)}
       \sim(nd)^{-e}
\]
for fixed \(e\).

For each fixed \(j\), ordered \(j\)-tuples of defects on disjoint vertex
sets give the leading contribution to \(\EE(Z)_j\).
Compatible overlapping tuples contribute \(o(1)\).
To see this, regard their prescribed pairs as a multigraph with \(v\)
vertices and \(e\) edges. Every vertex has degree at least two. A
connected component all of whose degrees equal two can only be a
single prescribed loop or a single prescribed double edge.
Thus a genuinely overlapping component has \(v<e\). There are
\(O(n^v)\) choices and each has probability \(O(n^{-e})\).
There are only finitely many overlap patterns for fixed \(d,j\).
Consequently
\[
 \EE(Z)_j\longrightarrow
 \left(\frac{d-1}{2}+\frac{(d-1)^2}{4}\right)^j.
\]
The finite inclusion-exclusion bounds for \(\PP(Z=0)\), followed by
increasing their order, give the asserted limit. Absence of these
defects is exactly simplicity.

For a triangle through a fixed vertex, summing over the other two
vertices and their stub choices gives
\(O(n^2)\,O(n^{-3})=O(n^{-1})\) in the matching model.
Conditioning on simplicity changes this bound by a bounded factor.
Finally, in any simple \(d\)-regular graph a uniform endpoint is a
uniform vertex. Each previously chosen endpoint has only \(O_d(1)\)
vertices equal or adjacent to it. A union bound over the fixed number
of sampled endpoints proves the last assertion.
\end{proof}

\subsection{The imported spectral estimates}

Put
\[
 S=\exp(\sqrt{\log n}).
\]
The following deliberately weak formulation is sufficient for our
argument.

\begin{proposition}[Spectral input]\label{prop:spectral-input}
There are constants \(c_0,c_1>0\), \(B<\infty\), and a fixed open interval
\(I_E\) containing \(E\), with closure contained in
\((-2\sqrt{d-1},2\sqrt{d-1})\), such that the following hold with probability
tending to one.

\begin{enumerate}
\item After omitting the eigenvalue associated with the constant vector,
the ordered eigenvalues are within \(n^{-c_0}\) of their
Kesten--McKay classical locations. In the normalization above, every
eigenvector entry is bounded by a deterministic \(n^{o(1)}\).
\item Uniformly for \(\Re\zeta\in I_E\) and \(S/n\le\Im\zeta\le1\),
\[
 \left|\frac1n\Tr(A-\zeta)^{-1}-m_d(\zeta)\right|
 \le(\log n)^B\bigl(n^{-c_1}+(n\Im\zeta)^{-c_1}\bigr).
\]
Except at at most \(n^{1-c_1}\) centers, the resolvent entries on a
center and its \(d\) neighbors satisfy the same error bound relative
to the corresponding infinite \(d\)-regular tree entries.
\end{enumerate}
These conclusions hold along all admissible \(n\). In particular,
\[
 \#\{i:|x_i|\le S^2\}=O(S^2)
\]
on an event of probability tending to one.
\end{proposition}

\begin{proof}
We use the fixed-degree results of Huang and Yau
\cite[Theorems 1.2, 1.4, and 4.2]{HuangYau}, in the precise version
arXiv:2102.00963v3.
Their matrix is \(A/\sqrt{d-1}\); hence their bulk coordinate is
\(E/\sqrt{d-1}\in(-2,2)\).
Rescaling the spectral parameter and the resolvent by this fixed factor
gives our normalization.

For clarity, the relevant error and geometry can be checked directly
in their Definition 4.1 and Equations (4.6)--(4.7).
Choose fixed parameters
\[
 0<c<1,\qquad a_{\mathrm{HY}}\ge12,\qquad
 b_{\mathrm{HY}}\ge25a_{\mathrm{HY}},\qquad
 0<\mathfrak r<c/32,
\]
and neighborhood radius
\(r_{\mathrm{HY}}=\mathfrak r\log_{d-1}n\), rounded to an integer.
Their Definition 1.1 and Proposition 2.1 give a tree-like event with
probability tending to one, on which all but a polynomially small
fraction of vertices have a tree neighborhood of radius
\((c/4)\log_{d-1}n\).
This is more than enough to contain every neighborhood needed for the
resolvent approximation between two vertices of the same star.
The tree extension with the semicircle weight in their Equation (2.12)
is then the ordinary regular-tree Green function.

In a fixed bulk interval, their error is bounded by a logarithmic
factor times
\[
 (d-1)^{-r_{\mathrm{HY}}}
 +(n\Im\zeta)^{-1/2}
 +(n\Im\zeta)^{-2/3}.
\]
The factor \(S\) dominates every fixed power of \(\log n\).
After decreasing \(c_1\) and increasing \(B\), this gives the stated
bounds and exceptional-set size. The same source provides polynomial
rigidity and polylogarithmic complete delocalization.
Its model and the quoted statements have no even-size restriction.

For the final assertion, apply the trace estimate at
\(\zeta=E+\mathrm iS^2/n\). Each eigenvalue with
\(|\lambda_i-E|\le S^2/n\) contributes at least
\(n/(2S^2)\) to \(\Im(\lambda_i-\zeta)^{-1}\).
The normalized imaginary trace is bounded, so the count is \(O(S^2)\).
\end{proof}

\begin{remark}
The estimates in Proposition~\ref{prop:spectral-input} are statements
about events in the unconditioned simple-graph ensemble. We never
condition the conclusion of Theorem~\ref{thm:main} on connectedness,
nonbipartiteness, or a tree-neighborhood property.
The bounded Laplace observables are unaffected in the limit by
discarding events of vanishing probability.
\end{remark}

\subsection{Covariance and logarithmic potential on the tree}

The covariance below is the Gaussian-wave covariance on the regular
tree; see \cite{ElonEigenvectors,CsokaEtAl}. We derive it together with
the logarithmic potential, so that the entropy normalization is explicit.

\begin{lemma}\label{lem:tree-identities}
Let \(q(\zeta)\) be the solution of
\[
 \zeta=q^{-1}+(d-1)q
\]
which tends to zero at infinity in the upper half-plane.
Then
\[
 -m_d(\zeta)=\frac1{\zeta-dq(\zeta)}.
\]
The regular-tree resolvent has diagonal entry \(m_d(\zeta)\) and
adjacent entry \((1+\zeta m_d(\zeta))/d\).
At \(\zeta=E+\mathrm i\eta\), its imaginary-part matrix on a star,
divided by \(\pi\rho\), tends at a polynomial rate in \(\eta\) to the
covariance matrix with variances one, center-neighbor covariance \(p\),
and distinct-neighbor covariance
\[
 \frac{E^2-d}{d(d-1)}.
\]
Furthermore,
\[
 k_0=\frac{1+Em_*}{d}
 \quad\Longrightarrow\quad
 k_0^2-k_0-m_*^2=0.
\]
For \(U(E)=\int\log|t-E|\rho_d(t)\dd t\),
\begin{equation}\label{eq:tree-potential}
 -U(E)=-\frac12\log d+\frac{d-1}{2}\log a
                       -\frac d4\log(1-p^2).
\end{equation}
Replacing \(E\) by \(E+\mathrm i\eta\) in the logarithmic potential
changes its real part by \(O(\eta)\).
\end{lemma}

\begin{proof}
A Schur complement at the root of a forward tree gives
\(q=(\zeta-(d-1)q)^{-1}\) in the convention \((\zeta-A)^{-1}\).
Attaching \(d\) forward branches yields the displayed formula for
\(-m_d\). The boundary value
\(\pi^{-1}\Im m_d(t+\mathrm i0)\) is the Kesten--McKay density above;
the branch has the required mass at infinity and no additional pole.

The resolvent equation at the center and symmetry among its neighbors
give the adjacent entry. Applying the same equation to a neighbor
column gives, for the common distinct-neighbor entry \(g_2\),
\[
 m_d(\zeta)+(d-1)g_2
       =\zeta\,\frac{1+\zeta m_d(\zeta)}d.
\]
Taking boundary imaginary parts proves the covariance formulas.
The tree formula is analytic across each side of a sufficiently small
bulk neighborhood, so these limits have a polynomial rate.
Substituting its boundary value into \(k_0=(1+Em_*)/d\) gives the
quadratic identity.

To compute the potential, integrate its analytic derivative
\(-m_d(\zeta)=q/(1-q^2)\) after changing from \(\zeta\) to \(q\).
Matching at infinity gives
\[
 \Re\int\log(\zeta-t)\rho_d(t)\dd t
       =-\log|q|-\frac{d-2}{2}\log|1-q^2|.
\]
At the bulk boundary,
\[
 |q|=(d-1)^{-1/2},\qquad
 |1-q^2|=\frac{\sqrt{d^2-E^2}}{d-1}.
\]
These identities give \eqref{eq:tree-potential} after substituting
\(a=d(1-p^2)/(d-1)\).
The boundary logarithmic integral exists because \(\rho_d\) is bounded
near \(E\). Finally,
\[
 \int \left(\log|t-E-\mathrm i\eta|-\log|t-E|\right)\rho_d(t)\dd t
 =\frac12\int\log\!\left(1+\frac{\eta^2}{(t-E)^2}\right)\rho_d(t)\dd t
 =O(\eta),
\]
by boundedness of the density near \(E\), the substitution
\(t-E=\eta u\), and the integrable tail of \(\log(1+u^{-2})\).
\end{proof}
